% !TEX program = lualatex
% Transparencias de Fundamentos de las Matemáticas I
% Clase 10: Relaciones binarias

\documentclass[aspectratio=169,11pt]{beamer}

\usepackage{fontspec}
\usepackage[spanish,es-nodecimaldot]{babel}
\usepackage{amsmath,amssymb,mathtools}
\usepackage{booktabs}
\usepackage{csquotes}
\usepackage{xcolor}

\usetheme{Madrid}
\usecolortheme{default}
\setbeamertemplate{navigation symbols}{}
\setbeamertemplate{footline}[frame number]

\definecolor{FdMBlue}{RGB}{20,55,100}
\definecolor{FdMRed}{RGB}{130,30,30}
\definecolor{FdMGreen}{RGB}{25,100,60}

\setbeamercolor{structure}{fg=FdMBlue}
\setbeamercolor{title}{fg=white,bg=FdMBlue}
\setbeamercolor{frametitle}{fg=white,bg=FdMBlue}
\setbeamercolor{block title}{fg=white,bg=FdMBlue}
\setbeamercolor{block body}{bg=blue!3}
\setbeamercolor{alerted text}{fg=FdMRed}

\setmainfont{Latin Modern Roman}
\setsansfont{Latin Modern Sans}
\setmonofont{Latin Modern Mono}

\newcommand{\N}{\mathbb{N}}
\newcommand{\Z}{\mathbb{Z}}
\newcommand{\Q}{\mathbb{Q}}
\newcommand{\R}{\mathbb{R}}
\newcommand{\C}{\mathbb{C}}
\newcommand{\Pow}{\mathcal{P}}
\newcommand{\id}{\operatorname{id}}
\newcommand{\mcd}{\operatorname{mcd}}
\newcommand{\mcm}{\operatorname{mcm}}

\title[FdM I -- Clase 10]{Relaciones binarias}
\subtitle{Definición, propiedades y ejemplos}
\author{Antonio Falcó Montesinos}
\institute{Universidad CEU Cardenal Herrera}
\date{Fundamentos de las Matemáticas I}

\begin{document}

\begin{frame}
  \titlepage
\end{frame}

\begin{frame}{Objetivos de la clase}
\begin{itemize}
  \item Definir relaciones binarias como subconjuntos de productos cartesianos.
  \item Identificar propiedades reflexiva, simétrica, antisimétrica y transitiva.
  \item Preparar relaciones de equivalencia y de orden.
\end{itemize}
\end{frame}

\begin{frame}{Mapa de la clase}
\tableofcontents
\end{frame}

\section{Definición}

\begin{frame}{Definición}
\begin{block}{Relación binaria}
\begin{itemize}
  \item Una relación binaria sobre \(A\) es un subconjunto de \(A\times A\).
  \item Si \((a,b)\in R\), se escribe a veces \(aRb\).
  \item La relación indica qué pares están conectados por la propiedad definida.
\end{itemize}
\end{block}
\end{frame}

\begin{frame}{Definición}
\begin{block}{Ejemplos}
\begin{itemize}
  \item En \(\mathbb{Z}\), la relación \(a\leq b\).
  \item En un conjunto de estudiantes, ``tener el mismo grupo de seminario''.
  \item En \(\mathbb{N}\), la relación de divisibilidad.
\end{itemize}
\end{block}
\end{frame}

\section{Propiedades}

\begin{frame}{Propiedades}
\begin{block}{Reflexiva y simétrica}
\begin{itemize}
  \item Reflexiva: para todo \(a\in A\), se cumple \(aRa\).
  \item Simétrica: si \(aRb\), entonces \(bRa\).
  \item Estas propiedades se verifican usando la definición de la relación.
\end{itemize}
\end{block}
\end{frame}

\begin{frame}{Propiedades}
\begin{block}{Antisimétrica y transitiva}
\begin{itemize}
  \item Antisimétrica: si \(aRb\) y \(bRa\), entonces \(a=b\).
  \item Transitiva: si \(aRb\) y \(bRc\), entonces \(aRc\).
  \item No hay que confundir simetría con antisimetría.
\end{itemize}
\end{block}
\end{frame}

\begin{frame}{Propiedades}
\begin{block}{Método de análisis}
\begin{itemize}
  \item Fijar el conjunto sobre el que está definida la relación.
  \item Traducir \(aRb\) usando la definición concreta.
  \item Comprobar una propiedad cada vez.
\end{itemize}
\end{block}
\end{frame}

\section{Profundización}

\begin{frame}{Profundización}
\begin{block}{Relación como subconjunto}
\begin{itemize}
  \item Definir una relación equivale a decir qué pares pertenecen a ella.
  \item Si \(R\subseteq A\times A\), entonces solo se relacionan elementos de \(A\).
  \item Cambiar el conjunto base puede cambiar las propiedades de la relación.
\end{itemize}
\end{block}
\end{frame}

\begin{frame}{Profundización}
\begin{block}{Ejemplo: divisibilidad}
\begin{itemize}
  \item En \(\mathbb{N}^{\ast}\), la relación \(a\mid b\) es reflexiva.
  \item Es transitiva: si \(a\mid b\) y \(b\mid c\), entonces \(a\mid c\).
  \item No es simétrica en general: \(2\mid 4\), pero \(4\nmid 2\).
\end{itemize}
\end{block}
\end{frame}

\begin{frame}{Profundización}
\begin{block}{Diagnóstico de propiedades}
\begin{itemize}
  \item Reflexiva: mirar pares \((a,a)\).
  \item Simétrica: intercambiar el orden del par.
  \item Transitiva: encadenar dos relaciones consecutivas.
\end{itemize}
\end{block}
\end{frame}

\begin{frame}{Cierre}
\begin{itemize}
  \item Una relación es un conjunto de pares.
  \item Las propiedades dependen de la definición concreta.
  \item Las dos familias clave serán equivalencias y órdenes.
\end{itemize}
\end{frame}

\begin{frame}{Trabajo recomendado}
\begin{itemize}
  \item Releer las secciones correspondientes del manual.
  \item Identificar definiciones, símbolos nuevos y enunciados principales.
  \item Preparar dudas y ejemplos para el seminario asociado.
\end{itemize}
\end{frame}

\end{document}
